\section{Cote 136: what ``\'egalise'' can mean homotopically}
\label{sec:136homotopy}

\paragraph{Archive and inherited inputs.}
This section concerns cote 136, physical p.~121, footer 120,
\texttt{ARCHIVE\_136\_scan121.tex}, with admissibility as on physical
pp.~117 and 119, footers 116 and 118, scans 117 and 119.  We use the
category $\mathcal A_k$ fixed in ALPOGE\_BRIDGES checkpoint 09,
AQ-136-COMPARISON, AQ-136-COTENSOR and AQ-136-COSPAN.  In particular
all bounds in the two admissibility conditions are uniform in the
simplex dimension.  The strict reading was refuted by AQ-136-STRICT;
we do not alter that conclusion.

Write
\[
 (P_J C)_n=C(\Delta[n]\times J)
\]
for the cotensor by a finite simplicial edge path $J$.  The underlying
graph of $J$ is an interval with at least one edge, but each edge may
have either orientation.  Its two distinct end vertices are denoted
$0$ and $1$.  A map into $P_J C$ is always required to preserve the
coefficient algebra $S_k=\Gamma_k(T)$, both gradings, differential,
products and divided powers.  Thus the coefficient paths are constant.

\paragraph{ARCH-136-HOM-DEFINITION (specified homotopy reading).}
For two strict arrows $f,g:A\rightrightarrows C$, write $f\simeq_s g$
if there is such a finite $J$ and a strict map
$H:A\to P_J C$ with endpoints $f,g$.  Equivalently this is the
equivalence relation generated by the ordinary simplicial cotensor
homotopies $A\to P_{\Delta[1]}C$.  This is a congruence on arrows:
precomposition acts on $H$, and postcomposition acts pointwise on
the cotensor.  Concatenation glues sections at their common endpoint;
reversal changes the order in which the oriented edges are traversed.
The constant homotopy uses a one-edge interval.

The resulting reading of (b) is
\[
 (b_s)\qquad
 \forall f,g:C\rightrightarrows D\quad
 \exists E\in\mathcal A_k,\ h:E\to C:
                         fh\simeq_s gh.
\]
This is an explicitly chosen reading, not a claim that the quotation
marks alone determine Grothendieck's intended definition.  It is
intrinsic to his simplicial generalized de Rham functors on physical
p.~111 (footer 110, scan 111), uses the very same algebraic structures,
and fixes the divided-power coefficient envelope.  It supplies a
precise interpretation of the all-arrow groupoid asked for on p.~121.
It is not the addition of an unweighted polynomial de Rham interval.

\paragraph{ARCH-136-PATHS (finite path cotensors and pullbacks).}
For every commutative ring $k$, $C\in\mathcal A_k$, and finite oriented
edge path $J$ as above, $P_J C$ is admissible.  Given a cospan
$A\xrightarrow{f} C\xleftarrow{g}B$, the strict algebraic limit
\[
 P_J(A,C,B)=(A\times B)\mathbin{\times}_{C\times C}P_J C
\]
is admissible.  Its projections to $A,B$ have composites to $C$
which are $\simeq_s$-homotopic.  No assertion that $A\times B$ itself
is admissible is involved.

\begin{proof}
The cotensor and the displayed limit inherit all operations pointwise.
Each leaf edge of $J$ can be collapsed to its adjacent vertex.  On an
edge with that vertex as source or target, the maximum or minimum map
of $\Delta[1]\times\Delta[1]$ provides a simplicial homotopy between
the identity and this retraction; it is the identity on the rest of
the path.  These formulas agree on the gluing vertex.  Successively
collapsing leaves proves that for any simplicial module $F$, $F^J$
has the same normalized homology as $F$.  The homotopies are linear
because the operations on sections are pointwise.  Consequently the
horizontal cutoff for $P_J C$ can be taken equal to that for $C$.

The simplicial set $\Delta[n]\times J$ is contractible.  The actual
comparison maps in AQ-136-COMPARISON therefore give its vertical
cohomology, in any fixed finite degree range and for a sufficiently
large weight, as $S_k^w$ in degree zero and zero in positive degrees.
The comparison bound is independent of the simplicial set.  In
particular it is uniform in $n$ (and can even be chosen independently
of the finite path).  This proves vertical admissibility of $P_J C$.

For completeness, endpoint surjectivity has the same uniformity.
If the normalized complex of $F=C_*^{w,q}$ is acyclic, a simplicial
linear map from
$k[\Delta[n]\times\{0,1\}]$ to $F$ extends to
$k[\Delta[n]\times J]$.  The normalized relative chains are free
on relative nondegenerate simplices, and the inclusion is split in
each degree.  Extend a chain map in increasing degree.  An unassigned
vertex can be assigned zero; the prescribed image of the boundary
of each higher cell is a cycle and has a primitive by acyclicity,
including $H_0=0$.  Dold--Kan converts this to the required extension.
Thus $(e_0,e_1):P_J C_n^{w,q}\to(C_n^{w,q})^2$ is surjective.

It follows, in the eventual range for each $q$, that the defining
limit has the short exact sequence of underlying simplicial modules
\[
 0\longrightarrow P_J(A,C,B)^{w,q}
 \longrightarrow A^{w,q}\oplus B^{w,q}\oplus(P_J C)^{w,q}
 \xrightarrow{(f-e_0,\,g-e_1)} (C^{w,q})^{\oplus2}
 \longrightarrow0.
\]
Exactness of normalization gives its horizontal acyclicity once
the finitely many rows on the right are acyclic.  For a vertical
cohomological range through $d$, choose one weight bound giving this
exactness through degree $d+1$ and the coefficient comparison for
all the finitely many terms through that range.  These bounds do
not depend on $n$.  The image of the displayed difference map agrees
with its target through degree $d+1$, so its cohomology through $d$
agrees with the target.  The map in degree-zero cohomology is
\[
  (S_k^w)^3\longrightarrow(S_k^w)^2,
  \qquad(a,b,t)\longmapsto(a-t,b-t).
\]
It is onto with diagonal kernel $S_k^w$, and the positive vertical
cohomology vanishes in the chosen range.  The long exact sequence
therefore proves precisely the required vertical admissibility.
The path component of the limit supplies the asserted homotopy.
\end{proof}

\paragraph{ARCH-136-FRACTIONS (right fractions after simplicial homotopy).}
Let $\mathcal H_k=\mathcal A_k/{\simeq_s}$.  All arrows of
$\mathcal H_k$ form a right multiplicative system.  The groupoid
obtained by inverting them is canonically the all-arrow groupoid
$\mathcal G_k$ of $\mathcal A_k$.  For strict parallel arrows one has
the exact criterion
\[
 f=g\text{ in }\mathcal G_k
 \quad\Longleftrightarrow\quad
 \exists E\in\mathcal A_k,\ h:E\to\operatorname{dom}(f),
                         \ fh\simeq_s gh.                 \tag{136.HF}
\]

\begin{proof}
Identities and composition cause no issue.  The right Ore square
for a cospan is given by AQ-136-COSPAN, whose square commutes in
$\mathcal H_k$.

Here is the cancellation axiom, which must not be inferred merely
from the existence of those squares.  Suppose $u:Y\to Z$ and
$f,g:X\rightrightarrows Y$ satisfy $uf\simeq_s ug$.  Choose a single
finite path $J$ encoding this entire homotopy, and form
\[
 E=(Y\times Y)\mathbin{\times}_{Z\times Z}P_JZ.
\]
By ARCH-136-PATHS it is admissible.  The homotopy and the two maps
define a strict $a:X\to E$.  Constant paths define the strict
diagonal $\delta:Y\to E$.  Apply AQ-136-COSPAN to
$X\xrightarrow a E\xleftarrow\delta Y$; this supplies an admissible
$R$ with strict $h:R\to X$, $v:R\to Y$, and $ah\simeq_s\delta v$.
Projection to the two $Y$ factors yields
$fh\simeq_s v$ and $gh\simeq_s v$, hence $fh\simeq_s gh$.
This proves cancellation for any finite chain of homotopies, even
when its intermediate maps to $Z$ do not factor through $u$.

These are exactly the right multiplicative-system axioms in
Stacks Project, Definition 4.27.1 (Tag 04VB).  Lemmas 4.27.11 and
4.27.16 construct the localization; Lemma 4.27.14, applied to the
common denominator $1_X$, gives (136.HF) for its morphisms.

To identify this groupoid with $\mathcal G_k$, any functor inverting
every strict arrow identifies the two endpoint maps of $P_JC$.
Both have the constant-path map $C\to P_JC$ as right inverse;
after inversion both are its inverse.  It consequently identifies
$\simeq_s$-homotopic maps, and factors through $\mathcal H_k$.
The two universal properties now agree.  All arrows of the
quotient category have strict representatives, so the denominator
$h$ furnished by the fraction criterion is an actual admissible
strict source arrow.  Conversely, a homotopy after $h$ becomes
an equality after inversion, and $h$ can be cancelled there.
\end{proof}

\paragraph{ARCH-136-HOM-SEMILOCAL (the requested implication from terminality).}
Over every commutative $k$, condition $(b_s)$ is equivalent to every
connected component of $\mathcal G_k$ being equivalent to a point.
In particular it follows from terminality of $\mathcal G_k$.
Together with common-source condition (a), it is equivalent to
terminality.  Hence $(b_s)$ holds over every ring in which all but
finitely many rational primes are invertible, including every
$\mathbb Q$-algebra, by AQ-136-SEMILOCAL and AQ-136-RATIONAL.

\begin{proof}
By (136.HF), $(b_s)$ says precisely that all original parallel
arrows have equal images in $\mathcal G_k$.  Within a connected
component, AQ-136-CONNECTED-SOURCES supplies one strict source for
any finite collection of objects.  Given two finite zigzags with
the same endpoints, choose such a source for all their vertices.
For every original edge, its composite with the chosen source map
and the chosen source map to its endpoint are parallel arrows;
their localized images are equal by $(b_s)$.  Both zigzags therefore
telescope to the same ratio of source arrows.  Thus each component
has exactly one morphism between every pair of objects.  The
converse follows immediately from (136.HF).  Common-source
condition (a) is connectedness by the cited inherited criterion;
nonemptiness holds by AQ-136-DP-MODEL.  The last claim uses the
supplied semilocal and rational terminality theorems.
\end{proof}

\paragraph{Integral state and geometric meaning.}
These results decide what follows from terminality without assuming
a strict calculus of fractions, which AQ-136-STRICT rules out.
Over $\mathbb Z$, they identify $(b_s)$ with absence of automorphisms
in the individual components, but do not prove that condition or
connectedness.  In Grothendieck's language, the comparison is between
admissible generalized de Rham theories over the fixed divided-power
envelope $\Gamma_k(T)$, by paths that preserve all divided powers.
It concerns the ordinary groupoid of these resolutions.  It neither
identifies this whole category with crystals on a particular scheme
nor asserts integral crystalline comparison without the hypotheses
of the relevant crystalline Poincare lemma.
